% ============================================================
% REZUMAT
% ============================================================

Această lucrare tratează proiectarea unui sistem e-medicine pentru analiza datelor
epidemiologice prin modele compartimentale, cu modelul SIDARTHE folosit ca
direcție clinică principală de explorare. Spre deosebire de modelele
SIR/SIS/SEIR/SIRS, care agregă dinamica în câteva stări canonice, SIDARTHE separă
infecțiile nedepistate, cazurile diagnosticate, formele simptomatice, formele
severe, vindecările și decesele.
Această granularitate face modelele mai potrivite pentru discutarea efectelor
testării, diagnosticării și severității clinice asupra evoluției unei epidemii.

Modelele SIR, SIS, SEIR și SIRS rămân prezente în lucrare ca repere
secundare, folosite pentru comparații și pentru verificarea comportamentului
numeric al implementărilor. Sistemul construit separă explicit stocarea datelor,
expunerea API-ului și execuția simulărilor.
Datele epidemiologice sunt persistate într-o schemă relațională compatibilă
PostgreSQL, rulată peste YugabyteDB, iar simulările sunt executate asincron de
runner-e specializate, coordonate prin Redis.

Lucrarea introduce un cadru explicit pentru compararea modelelor și calibrarea
parametrilor. Modelele sunt comparate structural, prin compartimente și ipoteze,
și numeric, prin operatori de observație, MAE, RMSE și nRMSE. Calibrarea este
formulată ca problemă de optimizare pe parametrii fiecărui model, cu accent pe
separarea dintre fereastra de antrenare și fereastra de test. În plus, este
integrată o extensie predictivă bazată pe LSTM, cu runner dedicat, folosită ca
model de serie temporală și ca posibilă corecție a reziduurilor modelelor
compartimentale.

Componenta web, alcătuită dintr-un API Go și două interfețe Vue.js separate,
oferă fluxurile necesare organizării datelor pe studii epidemiologice,
gestionării rolurilor de acces, încărcării seturilor de date, lansării
simulărilor, urmăririi statusului și vizualizării traiectoriilor
compartimentale. Aplicația principală este destinată cercetătorilor și
colaboratorilor dintr-un studiu, iar consola administrativă este rezervată
utilizatorilor cu drepturi globale de super-admin. Schema datelor
epidemiologice acoperă, dincolo de cele patru câmpuri SIR-centrice, nouă
metrici suplimentare (decese, spitalizări, terapie intensivă, vaccinări,
teste, indicele Oxford de stringență, $R_t$ observat) propagate din sursa
OWID; aceste valori nu sunt folosite de runner-ele compartimentale built-in,
dar sunt expuse integral scripturilor scrise de cercetător într-un editor
web dedicat, ceea ce permite formularea de modele ad-hoc fără modificări
de cod în API. Întregul sistem este
deployabil pe un cluster Kubernetes prin manifeste declarative, pe o platformă
cloud privată cu factor de replicare~3 pentru baza de date. Sistemul este validat empiric printr-o campanie de benchmark executată prin
aplicație (180 de rulări ODE cu cinci modele compartimentale, trei limbaje
și patru orizonturi temporale) și prin calibrarea SIR pe valul 1 al
pandemiei COVID-19 în Romania, Italia, Germania și Spania, cu analiza
identifiabilității practice prin scanarea peisajului de pierdere. Accentul
lucrării cade pe legătura dintre modelul matematic, implementarea numerică
și arhitectura software care face posibilă evaluarea reproductibilă a
rezultatelor pentru modele cu grade diferite de complexitate.

\bigskip
\noindent\textbf{Cuvinte cheie:} model SIDARTHE, modele SIR, epidemiologie computațională, e-medicine,
calibrare de parametri, LSTM, RBAC, YugabyteDB, Kubernetes, Go, Vue.js,
benchmark de performanță.
